% GENERATED FILE. Edit canonical lessons, then run npm run generate:books.
% canonical-source-hash: fnv1a64:bbf049d3cbd74ced
% canonical-lessons: LA-C98-dives, LA-C98-pauper, LA-C98-stultus, LA-C98-satis, LA-C98-nimis

\chapter{Rich, Poor, Foolish, Enough, Too Much}
\label{ch:la-rich-poor-foolish-enough-too-much}
\hlchaptermodality{98}

\begin{chapteropening}
\textbf{By the end of this chapter:} \emph{I can say rich, poor and foolish, and enough and too much.}

Complete the last lesson of chapter 98: I can say rich, poor and foolish, and enough and too much.
\end{chapteropening}

\section[dīves, dīvitis]{dīves, dīvitis — rich}

\label{lesson:LA-C98-dives}



\begin{quote}
Before the new one: say the Latin for brave, then the Latin for wise.
\end{quote}

\subsection*{You'll want to know: dīves, dīvitis}
\textbf{dīves, dīvitis} — \textquotedblleft{}rich\textquotedblright{}.

\begin{cousinweb}
\textbf{dīves} is often derived from \emph{dīvus}, \textquotedblleft{}godlike, divine\textquotedblright{} (Lewis \& Short): the rich as those the gods favour. That account is old and still uncertain.
\end{cousinweb}

\begin{grammarlens}[title={What you've built}]
One more word for this chapter.
\end{grammarlens}

\subsection*{Your turn}
\begin{itemize}
\raggedright
  \item \emph{Say it:} \emph{dīves}
  \item \emph{Say it:} \emph{dīves}, once more
  \item \emph{Say it:} say \emph{dīves}
  \item \emph{Recall:} say the Latin for brave, then the Latin for wise, then say \emph{dīves} again
\end{itemize}

\begin{tcolorbox}[breakable,colback=teal!4,colframe=teal!35!black,title={Before you move on}]
What does dīves, dīvitis mean? (\textquotedblleft{}Rich\textquotedblright{}.) Say it once more.
\end{tcolorbox}
\section[pauper, pauperis]{pauper, pauperis — poor}

\label{lesson:LA-C98-pauper}



\begin{quote}
Before the new one: say the Latin for wise, then the Latin for rich.
\end{quote}

\subsection*{You'll want to know: pauper, pauperis}
\textbf{pauper, pauperis} — \textquotedblleft{}poor\textquotedblright{}. English \emph{pauper} and \emph{poverty}.

\begin{cousinweb}
\textbf{pauper} is usually explained as \emph{pau-}, \textquotedblleft{}little\textquotedblright{}, as in \emph{paucus}, \textquotedblleft{}few\textquotedblright{}, plus the root of \emph{pariō}, \textquotedblleft{}to produce\textquotedblright{} (Lewis \& Short): one who gets or produces little.
\end{cousinweb}

\begin{grammarlens}[title={What you've built}]
One more word for this chapter.
\end{grammarlens}

\subsection*{Your turn}
\begin{itemize}
\raggedright
  \item \emph{Say it:} \emph{pauper}
  \item \emph{Say it:} \emph{pauper}, once more
  \item \emph{Say it:} say \emph{pauper}
  \item \emph{Recall:} say the Latin for wise, then the Latin for rich, then say \emph{pauper} again
\end{itemize}

\begin{tcolorbox}[breakable,colback=teal!4,colframe=teal!35!black,title={Before you move on}]
What does pauper, pauperis mean? (\textquotedblleft{}Poor\textquotedblright{}.) Say it once more.
\end{tcolorbox}
\section[stultus, stulta, stultum]{stultus, stulta, stultum — foolish}

\label{lesson:LA-C98-stultus}



\begin{quote}
Before the new one: say the Latin for rich, then the Latin for poor.
\end{quote}

\subsection*{You'll want to know: stultus, stulta, stultum}
\textbf{stultus, stulta, stultum} — \textquotedblleft{}foolish\textquotedblright{}.

\begin{cousinweb}
\textbf{stultus} is probably related to \emph{stolidus}, \textquotedblleft{}dull, stupid\textquotedblright{} (Lewis \& Short), and both are often traced to Proto-Indo-European \emph{*ster-}, \textquotedblleft{}stiff\textquotedblright{} (EtymDB 2.1): a fool as one stiff of mind.
\end{cousinweb}

\begin{grammarlens}[title={What you've built}]
One more word for this chapter.
\end{grammarlens}

\subsection*{Your turn}
\begin{itemize}
\raggedright
  \item \emph{Say it:} \emph{stultus}
  \item \emph{Say it:} \emph{stultus}, once more
  \item \emph{Say it:} say \emph{stultus}
  \item \emph{Recall:} say the Latin for rich, then the Latin for poor, then say \emph{stultus} again
\end{itemize}

\begin{tcolorbox}[breakable,colback=teal!4,colframe=teal!35!black,title={Before you move on}]
What does stultus, stulta, stultum mean? (\textquotedblleft{}Foolish\textquotedblright{}.) Say it once more.
\end{tcolorbox}
\section[satis]{satis — enough}

\label{lesson:LA-C98-satis}



\begin{quote}
Before the new one: say the Latin for poor, then the Latin for foolish.
\end{quote}

\subsection*{You'll want to know: satis}
\textbf{satis} — \textquotedblleft{}enough\textquotedblright{}. English \emph{satisfy}.

\begin{cousinweb}
\textbf{satis} is inherited, usually traced to Proto-Indo-European \emph{*sh2tis}, \textquotedblleft{}fullness\textquotedblright{}, from \emph{*seh2-}, \textquotedblleft{}to satisfy, to fill\textquotedblright{} (EtymDB 2.1). English \emph{sad} is a cousin: it first meant \textquotedblleft{}having had enough\textquotedblright{}.
\end{cousinweb}

\begin{grammarlens}[title={What you've built}]
One more word for this chapter.
\end{grammarlens}

\subsection*{Your turn}
\begin{itemize}
\raggedright
  \item \emph{Say it:} \emph{satis}
  \item \emph{Say it:} \emph{satis}, once more
  \item \emph{Say it:} say \emph{satis}
  \item \emph{Recall:} say the Latin for poor, then the Latin for foolish, then say \emph{satis} again
\end{itemize}

\begin{tcolorbox}[breakable,colback=teal!4,colframe=teal!35!black,title={Before you move on}]
What does satis mean? (\textquotedblleft{}Enough\textquotedblright{}.) Say it once more.
\end{tcolorbox}
\section[nimis]{nimis — too much}

\label{lesson:LA-C98-nimis}



\begin{quote}
Before the new one: say the Latin for foolish, then the Latin for enough.
\end{quote}

\subsection*{You'll want to know: nimis}
\textbf{nimis} — \textquotedblleft{}too much\textquotedblright{}.

\begin{cousinweb}
\textbf{nimis} probably begins with the negative \emph{ne-}, \textquotedblleft{}not\textquotedblright{}, but the rest is disputed: some link it to the stem of \emph{minus}, \textquotedblleft{}less\textquotedblright{}, while Lewis \& Short tie it to a root for measuring.
\end{cousinweb}

\begin{grammarlens}[title={What you've built}]
One more word for this chapter.
\end{grammarlens}

\subsection*{Your turn}
\begin{itemize}
\raggedright
  \item \emph{Say it:} \emph{nimis}
  \item \emph{Say it:} \emph{nimis}, once more
  \item \emph{Say it:} say \emph{nimis}
  \item \emph{Recall:} say the Latin for foolish, then the Latin for enough, then say \emph{nimis} again
\end{itemize}

\begin{tcolorbox}[breakable,colback=teal!4,colframe=teal!35!black,title={Before you move on}]
What does nimis mean? (\textquotedblleft{}Too much\textquotedblright{}.) Say it once more.
\end{tcolorbox}
